% ==============================================================================
% Phân mục 04 [TRỌNG TÂM NHIỆM VỤ 3.2]: Kiến trúc Phần mềm & Giao thức RESTful
% ==============================================================================
\section{Kiến Trúc Phần Mềm Đa Tầng và Giao Thức Hợp Đồng Dữ Liệu RESTful (Software Architecture)}
\label{sec:software_architecture}

Để chuyển hóa các mô hình học sâu đa phương thức từ môi trường thử nghiệm rời rạc thành một hệ thống hỗ trợ quyết định lâm sàng có khả năng vận hành thực tế tại bàn khám bệnh, M-CDSS được thiết kế theo kiến trúc vi dịch vụ 3 tầng tách biệt (\textit{Decoupled Three-Tier Microservice Architecture}). Toàn bộ kiến trúc được minh họa chi tiết tại Hình~\ref{fig:system_architecture}.

% ==============================================================================
% Hình 1: Sơ đồ Kiến trúc Phần mềm Đa tầng M-CDSS (Vector TikZ)
% ==============================================================================
\begin{figure}[htbp]
\centering
\resizebox{\textwidth}{!}{%
\begin{tikzpicture}[node distance=1.2cm and 1.6cm]

    % --- TẦNG 1: PRESENTATION TIER (NEXT.JS 14 FRONTEND - PORT 3000) ---
    \node[mcdssBox] (ui_form) {
        \textbf{PatientDataForm.tsx}\\
        \scriptsize Nhập 12 sinh hiệu \& bệnh án
    };
    \node[mcdssBox, below=0.7cm of ui_form] (ui_card) {
        \textbf{DiagnosticCard.tsx}\\
        \scriptsize Thanh xác suất hiệu chuẩn ($\text{ECE}<0.05$)
    };
    \node[mcdssBox, below=0.7cm of ui_card] (ui_ev) {
        \textbf{EvidenceViewer \& Drawer}\\
        \scriptsize Highlight $[start, end]$ \& CPGs
    };
    \node[mcdssBox, below=0.7cm of ui_ev] (ui_dca) {
        \textbf{DCACurveChart.tsx}\\
        \scriptsize Đồ thị Net Benefit ($p_t \in [0.05, 0.60]$)
    };

    \begin{scope}[on background layer]
        \node[mcdssGroup, fill=blue!3, fit=(ui_form) (ui_dca), label={[font=\footnotesize\bfseries\sffamily, text=mcdssPrimary]above:Tầng Giao Diện (Next.js 14 :3000)}] (tier_frontend) {};
    \end{scope}

    % --- TẦNG 2: APPLICATION & API GATEWAY TIER (FASTAPI - PORT 8000) ---
    \node[mcdssService, right=2.4cm of ui_form] (api_pred) {
        POST /api/v1/predict/mock\\
        \normalfont\scriptsize Multimodal \& Baseline Inference
    };
    \node[mcdssService, right=2.4cm of ui_card] (api_agent) {
        POST /api/v1/agent/explain\\
        \normalfont\scriptsize Multi-Agent RAG \& Grounding
    };
    \node[mcdssService, right=2.4cm of ui_ev] (api_eval) {
        GET /api/v1/evaluation/benchmark\\
        \normalfont\scriptsize Metrics, DCA Points \& PROBAST
    };
    \node[mcdssService, right=2.4cm of ui_dca] (api_health) {
        GET /health \& Pydantic V2\\
        \normalfont\scriptsize Kiểm soát biên sinh lý \& CORS
    };

    \begin{scope}[on background layer]
        \node[mcdssGroup, fill=teal!3, fit=(api_pred) (api_health), label={[font=\footnotesize\bfseries\sffamily, text=mcdssSecondary]above:Tầng Điều Phối RESTful (FastAPI :8000)}] (tier_api) {};
    \end{scope}

    % --- TẦNG 3: AI CORE & DATA PERSISTENCE TIER ---
    \node[mcdssBox, right=2.4cm of api_pred, text width=3.9cm] (core_fusion) {
        \textbf{Multimodal Fusion Engine}\\
        \scriptsize FT-Transformer + Bio\_ClinicalBERT\\
        Cross-Attention + Temp Scaling ($T$)
    };
    \node[mcdssAgent, right=2.4cm of api_agent, text width=3.9cm] (core_agent) {
        \textbf{Agentic RAG \& Grounder}\\
        \scriptsize LangGraph Orchestrator +\\
        Kiểm chứng vị trí ký tự $\mathcal{O}(N)$
    };
    \node[mcdssBox, right=2.4cm of api_eval, text width=3.9cm] (core_qdrant) {
        \textbf{Qdrant Vector DB (:6333)}\\
        \scriptsize Kho phác đồ GINA, GOLD, ESC\\
        Hybrid Search (Dense + BM25)
    };
    \node[mcdssBox, right=2.4cm of api_health, text width=3.9cm] (core_store) {
        \textbf{Strict Split Benchmarks}\\
        \scriptsize Tập kiểm thử độc lập ($N=302$)\\
        0\% Patient-level Leakage
    };

    \begin{scope}[on background layer]
        \node[mcdssGroup, fill=orange!3, fit=(core_fusion) (core_store), label={[font=\footnotesize\bfseries\sffamily, text=mcdssAccent]above:Tầng Lõi AI \& Lưu Trữ Tri Thức}] (tier_core) {};
    \end{scope}

    % --- KẾT NỐI MŨI TÊN (REST JSON & INTERNAL CALLS) ---
    \draw[mcdssArrow, <->] (ui_form.east) -- node[above, font=\tiny\sffamily] {JSON Payload} (api_pred.west);
    \draw[mcdssArrow, <->] (ui_card.east) -- node[above, font=\tiny\sffamily] {Calibrated $P(y|x)$} (api_pred.west);
    \draw[mcdssArrow, <->] (ui_ev.east) -- node[above, font=\tiny\sffamily] {Spans + CPGs} (api_agent.west);
    \draw[mcdssArrow, <->] (ui_dca.east) -- node[above, font=\tiny\sffamily] {DCA \& PROBAST} (api_eval.west);

    \draw[mcdssArrow] (api_pred.east) -- (core_fusion.west);
    \draw[mcdssArrow] (api_agent.east) -- (core_agent.west);
    \draw[mcdssDashedArrow, <->] (core_agent.south) -- node[right, font=\tiny\sffamily] {Top-$k$ Chunks} (core_qdrant.north);
    \draw[mcdssArrow] (api_eval.east) -- (core_store.west);

\end{tikzpicture}%
}
\caption{Kiến trúc phần mềm 3 tầng đóng gói trong mạng nội bộ Docker (\texttt{mcdss-internal-net}), kết nối giao diện bác sĩ (Next.js 14), cổng điều phối hợp đồng dữ liệu (FastAPI) và lõi suy luận đa phương thức cùng kho vector y văn (Qdrant).}
\label{fig:system_architecture}
\end{figure}


\subsection{Phân tách Trách nhiệm Ba Tầng Dịch vụ}
Kiến trúc hệ thống tuân thủ nguyên lý tách biệt mối quan tâm (\textit{Separation of Concerns}) thông qua ba dịch vụ độc lập giao tiếp qua mạng nội bộ ảo \texttt{mcdss-internal-net}:
\begin{enumerate}[leftmargin=1.8em, itemsep=0.25em]
    \item \textbf{Tầng Giao diện Lâm sàng (\textit{Presentation Tier} --- Next.js 14, Cổng 3000):} Được xây dựng bằng \texttt{React 18} và \texttt{TypeScript} theo kiến trúc \textit{App Router}. Tầng này cung cấp bàn làm việc hội chẩn tương tác cho bác sĩ (\texttt{PatientDataForm}, \texttt{DiagnosticCard}, \texttt{EvidenceViewer}, \texttt{GuidelineDrawer}) và bảng điều khiển thẩm định lâm sàng (\texttt{DCACurveChart}). Để tối ưu hóa dung lượng đóng gói, ứng dụng được biên dịch ở chế độ độc lập (\texttt{output: "standalone"}) qua quy trình Docker đa tầng (\textit{Multi-stage build}).
    \item \textbf{Tầng Điều phối \& Hợp đồng Dữ liệu (\textit{Application \& API Gateway Tier} --- FastAPI, Cổng 8000):} Đóng vai trò cổng giao tiếp RESTful trung tâm vận hành trên máy chủ bất đồng bộ \texttt{Uvicorn}. Tầng này thực thi kiểm định biên sinh lý chặt chẽ bằng \texttt{Pydantic V2} trước khi chuyển luồng dữ liệu tới lõi suy luận đa phương thức hoặc đồ thị tác tử giải thích.
    \item \textbf{Tầng Lõi Suy luận AI \& Cơ sở dữ liệu Vector (\textit{AI Core \& Persistence Tier} --- PyTorch \& Qdrant, Cổng 6333):} Bao gồm mô hình dung hợp \texttt{MultimodalEHRClassifier} đã hiệu chuẩn nhiệt độ và cơ sở dữ liệu vector \texttt{Qdrant} lưu trữ các phân đoạn phác đồ điều trị y khoa chính thống (GINA, GOLD, ESC, Bộ Y tế).
\end{enumerate}

\subsection{Thiết kế Giao thức RESTful và Kiểm định Biên Sinh lý Tự động}
Trong tin học y tế, việc truyền nhầm đơn vị đo lường hoặc nhập các giá trị sinh hiệu phi sinh lý (ví dụ: $\text{SpO}_2 > 100\%$ hoặc nhịp tim âm) có thể khiến mô hình học sâu đưa ra dự đoán sai lệch nguy hiểm. Do đó, mọi yêu cầu HTTP gửi tới tầng Backend đều bắt buộc đi qua lớp kiểm định hợp đồng dữ liệu \texttt{TabularFeatures} (\texttt{src/api/schemas.py}):
\begin{itemize}[leftmargin=1.8em, itemsep=0.2em]
    \item \textbf{Ràng buộc miền giá trị sinh lý (\textit{Physiological Domain Constraints}):} Tuổi bệnh nhân $\in [0, 120]$ năm; Huyết áp tâm thu $\in [40, 300]\text{ mmHg}$; Huyết áp tâm trương $\in [20, 200]\text{ mmHg}$; Nhịp tim $\in [20, 300]\text{ bpm}$; Nhịp thở $\in [4, 80]\text{ lần/phút}$; Thân nhiệt $\in [30.0, 45.0]^\circ\text{C}$; Độ bão hòa oxy máu $\text{SpO}_2 \in [50.0, 100.0]\%$.
    \item \textbf{Ràng buộc toàn vẹn văn bản bệnh án (\textit{Clinical Note Validation}):} Trường \texttt{clinical\_note} bắt buộc có độ dài tối thiểu $L_{\min} \ge 10$ ký tự sau khi cắt bỏ khoảng trắng thừa; mọi yêu cầu vi phạm bị chặn ngay tại tầng cổng với mã trạng thái \texttt{HTTP 422 Unprocessable Entity}.
    \item \textbf{Định chuẩn Tiên đề Xác suất (\textit{Probability Axiom Enforcement}):} Tại đầu ra, lớp kiểm định \texttt{PredictionResponse} xác nhận vector xác suất dự đoán trên 4 nhóm bệnh thỏa mãn $\sum_{c=1}^{4} \hat{p}_c = 1.0 \pm 10^{-3}$ trước khi trả về cho trình duyệt của bác sĩ.
\end{itemize}

Bảng~\ref{tab:rest_endpoints} tổng hợp toàn bộ hệ thống điểm cuối RESTful API phục vụ luồng hội chẩn và kiểm định lâm sàng.

\begin{table}[htbp]
\centering
\caption{Đặc tả các điểm cuối giao thức RESTful API trọng tâm của hệ thống M-CDSS.}
\label{tab:rest_endpoints}
\footnotesize
\begin{tabularx}{\textwidth}{@{}llX@{}}
\toprule
\textbf{Phương thức \& Đường dẫn} & \textbf{Schema Đầu ra Chính} & \textbf{Chức năng Lâm sàng \& Đảm bảo Kỹ thuật} \\
\midrule
\texttt{GET /health} & \texttt{HealthResponse} & Kiểm tra trạng thái sống (\textit{Liveness/Readiness}) của container API và kết nối nội bộ Docker. \\[0.4em]
\texttt{POST /api/v1/predict/mock} & \texttt{PredictionResponse} & Thực thi suy luận đa phương thức (hoặc đối chứng với 3 mô hình Baseline), trả về phân phối xác suất đã hiệu chuẩn ($\text{ECE} < 0.05$) và trọng số chú ý chéo. \\[0.4em]
\texttt{POST /api/v1/agent/explain} & \texttt{AgentExplainResponse} & Kích hoạt đồ thị hội chẩn đa tác tử \texttt{LangGraph}, truy xuất phác đồ GINA/GOLD và trả về danh sách câu bằng chứng đã xác thực tọa độ ký tự $[start, end]$. \\[0.4em]
\texttt{GET /api/v1/benchmarks/baselines} & \texttt{BaselineBenchmarkList} & Xuất ma trận chỉ số định lượng của 3 mô hình đơn phương thức trên tập kiểm thử độc lập ($N=302$). \\[0.4em]
\texttt{GET /api/v1/benchmarks/ablation} & \texttt{AblationBenchmarkList} & Cung cấp kết quả thực nghiệm cắt bỏ thành phần (\textit{Ablation Study}) chứng minh đóng góp của tầng Cross-Attention và Temperature Scaling. \\[0.4em]
\texttt{GET /api/v1/evaluation/benchmark} & \texttt{EvaluationBenchmark} & Cung cấp tọa độ đường cong quyết định DCA ($p_t \in [0.05, 0.60]$) và bảng kiểm 4 miền PROBAST (\textit{Low Risk of Bias}). \\
\bottomrule
\end{tabularx}
\end{table}

